%% ma: L11-B16-0037
%% lop: 11
%% chuong: 5
%% bai: 16
%% dang: 11.16.4
%% loai: TN4
%% muc: VD
%% dap_an: "A"
%% nguon: "tai-lieu-hoc-tap-mon-toan-lop-11-theo-chuong-trinh-giao-duc-pho-thong-2018.pdf (Lớp Toán Thầy Duy, Phạm Lê Duy, THPT Chu Văn An), Chương 3 Bài 2, Phần 1, Câu 25 (tr.366)"
%% kien_thuc: "Lớp 11 Bài 16 (giới hạn của hàm số)"
%% trang_thai: chinh-thuc
%% ghi_chu: "Đã duyệt 03/11/2026."
\begin{ex}
Cho $\lim\limits_{x\to+\infty}\left(\dfrac{x^2+3x+1}{x+1}+ax+b\right)=1$. Khi đó giá trị của biểu thức $T=a+b$ bằng
\choice{\True $-2$}{$0$}{$1$}{$2$}
\loigiai{\sloppy $\dfrac{x^2+3x+1}{x+1}=x+2-\dfrac1{x+1}$ nên biểu thức $=(1+a)x+(2+b)-\dfrac1{x+1}$. Để có giới hạn hữu hạn thì $1+a=0$, rồi $2+b=1$. Vậy $a=-1$, $b=-1$, $T=-2$.}
\end{ex}
